\documentclass[10pt,a4paper]{amsart}
\usepackage[margin=19mm]{geometry}
\usepackage{amsmath,amssymb,mathtools}
\usepackage[scr=rsfs,cal=euler]{mathalfa}
\usepackage{xfrac}
\usepackage[unicode,hidelinks,bookmarks=false]{hyperref}
\newcommand{\ie}{i.e.\ }
\newcommand{\cf}{cf.\ }
\newcommand{\resp}{respectively\ }
\newcommand{\Subsection}{\subsection}
\providecommand{\nobreakdash}{\nobreak\hbox{-}}
\setlength{\emergencystretch}{2em}
\multlinegap0pt
\usepackage{mdframed}
\usepackage{mdframed}
\usepackage{mdframed}
\usepackage{mdframed}
\usepackage{bm}
\usepackage{amssymb}
\usepackage{xcolor}
\usepackage{algorithm}\usepackage{algpseudocode}
\usepackage{mathtools}
\usepackage{tikz}
\usepackage{float}
\usepackage{mdframed}
\newtheorem{lemma}{Lemma}
\def \bmv {\bm{v}}
\def \bmw {\bm{w}}
\def \calA {\mathcal{A}}
\title{Frequency Domain Biot--Allard Equations for Isotropic and Anisotropic Poroelastic Media: Two-field formulations and iterative splitting}
\author{Morten Jakobsen, Jakob Seierstad Stokke, Kundan Kumar, Florin Adrian Radu}
\date{Source: arXiv:2609.07388v1 (CC BY 4.0)}
\begin{document}
\maketitle

\noindent\emph{Source and licence.} Frequency Domain Biot--Allard Equations for Isotropic and Anisotropic Poroelastic Media: Two-field formulations and iterative splitting. Version arXiv:2609.07388v1 (CC BY 4.0), distributed by arXiv under the Creative Commons Attribution 4.0 International licence, http://creativecommons.org/licenses/by/4.0/, as recorded in the arXiv licence metadata for this exact version and shown on the abstract page https://arxiv.org/abs/2609.07388v1. Full licence notice: \texttt{licenses/arxiv-2609.07388-CC-BY-4.0.txt}.\par\smallskip

\noindent\textit{Editorial scope.} Counted unit: the article's lemma lemma (source0.tex lines 1981--2005). The article's complete original proof of it is delivered with it (source0.tex lines 2007--2043, one \texttt{proof} environment of 37 lines). No mathematics is written, repaired or re-derived: the statement and the proof are the article's own lines, in the article's own notation, and no retained line is substituted or re-flowed.  No further source-local context is needed: every cross-reference of every retained line resolves inside the two delivered spans.  Nothing was omitted from any retained span, so the package carries no omission record.  No retained line contains a citation, so the wrapper carries no bibliography and no bibliography entry.  No retained line contains a figure environment, an image-inclusion command or drawing code, so no image is delivered and the package declares no asset dependency.  Editorial formatting only, in addition: an AMS \texttt{amsart} preamble that restates the article's own declarations and macros used by the retained lines, namely the environments \texttt{lemma} on separate counters, together with the standard packages those lines need.  The retained environments print in this document as Lemma 1 in order of appearance, whereas the article prints the same items with the numbering of its own full document.  The complete native source of the article as distributed in the arXiv e-print for this exact version is bundled unchanged as \texttt{sources/209624/source0.tex}.
	\begin{lemma}[Adjointness of the dynamic $B$--coupling terms]
		%Let $\Omega\subset\mathbb{R}^d$ be a bounded Lipschitz domain and assume homogeneous boundary conditions
		%(e.g.\ $H^1_0(\Omega)$ traces). 
		Let $\bmv\in \bm{H}_{0}^{1}$ and $\dot{p}\in H^{1}_{0}$. Consider the frequency-dependent dynamic coupling operators
		\[
		(A_{12}^B \dot p)_j := -\,i\omega \rho_f\, B_{j k}\,\partial_k \dot p,
		\qquad
		A_{21}^B \mathbf{v} := -\,i\omega \rho_f\, \partial_k\!\big(B_{k j}\,v_j\big),
		\]
		with a sufficiently regular tensor field $B(\omega,\mathbf{x})\in\mathbb{C}^{d\times d}$.
		%With respect to the complex $L^2$ inner products
		%\[
		%\langle \mathbf{a},\mathbf{b}\rangle_{\bm{L}^{2}} = \int_\Omega a_j\,\overline{b_j}\,d\bm{x},
		%\qquad
		%\langle \phi,\psi\rangle_{L^2(\Omega)} = \int_\Omega \phi\,\overline{\psi}\,d\bm{x},
		%\]
		Then the off--diagonal operators satisfy
		\[
		A_{12}^B \;=\; (A_{21}^B)^{*}
		\quad\Longleftrightarrow\quad
		B_{k j} = \overline{B_{j k}},
		\]
		with respect to the complex $L^2$ inner products
		In particular, the condition holds when $\bm{B}$ is hermitian.
	\end{lemma}

	\begin{proof}
		\emph{Step 1 (integration by parts).}
		For $\bmw\in \bm{H}_{0}^{1}$ and $\dot{p}\in H^{1}_{0}$, we have
		\[
		\begin{aligned}
			\langle \calA_{12}^B \dot p, \mathbf{w}\rangle
			&= \int_\Omega \big(-i\omega\rho_f\, B_{j k}\,\partial_k \dot p\big)\,\overline{w_j}\,d\bm{x}
			= -\,i\omega\rho_f \int_\Omega B_{j k}\,(\partial_k \dot p)\,\overline{w_j}\,d\bm{x}
			\\
			&= i\omega\rho_f \int_\Omega \dot p\,\partial_k\big(B_{j k}\,\overline{w_j}\big)\,d\bm{x}
			\qquad\text{(by IBP; boundary term vanishes)}
			\\
			&= i\omega\rho_f \int_\Omega \dot p\,\overline{\partial_k\big(\overline{B_{j k}}\,w_j\big)}\,d\bm{x}
			= \omega\rho_f \int_\Omega \dot p\,\overline{(-i)\partial_k\big(\overline{B_{j k}}\,w_j\big)}\,d\bm{x}.
		\end{aligned}
		\]
		
		\emph{Step 2 (adjoint identification).}
		By the definition of the adjoint,
		\[
		\langle A_{12}^B \dot p, \mathbf{w}\rangle
		= \langle \dot p,(A_{12}^B)^{*}\mathbf{w}\rangle
		\quad\Rightarrow\quad
		(A_{12}^B)^{*}\mathbf{w} = -i\omega\rho_f\,\partial_k\big(\overline{B_{j k}}\,w_j\big).
		\]
		On the other hand,
		\[
		A_{21}^B \mathbf{w} = -\,i\omega\rho_f\,\partial_k\big(B_{k j}\,w_j\big).
		\]
		Thus, $A_{12}^B = (A_{21}^B)^{*}$ iff
		\[
		\partial_k\big(B_{k j}\,w_j\big) = \partial_k\big(\overline{B_{j k}}\,w_j\big)\quad\forall\,\mathbf{w}
		\;\;\Longleftrightarrow\;\;
		B_{k j} = \overline{B_{j k}},
		\]
		which is the Hermiticity of $\bm{B}$.
	\end{proof}

\end{document}
